% =================================================
% 文件: Linux.tex
% 章节: Linux 操作系统
% 版本: v1.0
% =================================================

\documentclass[12pt,UTF8]{ctexbook}

\input{../common/page}

\xeCJKsetup{AutoFallBack=true}
\setCJKfamilyfont{hei}{SimHei}
\setCJKfamilyfont{kai}{KaiTi}

\usepackage{titletoc}
\titlecontents{chapter}[0pt]{\vspace{3mm}\bf\addvspace{2pt}\filright}{\contentspush{\thecontentslabel\hspace{0.8em}}}{}{\titlerule*[8pt]{.}\contentspage}

\usepackage[colorlinks,linkcolor=blue,citecolor=blue,urlcolor=blue]{hyperref}

\ctexset{
	part/name= {第,卷},
	part/number={\chinese{part}},
	chapter/name={第,章},
	chapter/number={\arabic{chapter}}
}

\usepackage{graphicx}
\graphicspath{{Images/}}

\usepackage[perpage]{footmisc}

\usepackage{enumitem}

\title{\heiti\zihao{0} Linux 操作系统入门指南}
\author{WangFei}
\date{\today}

\begin{document}

\maketitle
\tableofcontents

\frontmatter
\chapter{前言}

本指南面向初学者，详细介绍 Linux 操作系统的基础概念、命令行操作、文件系统、用户管理、进程管理、系统配置等知识。

\mainmatter

\chapter{Linux 基础概念}
\label{chap:linux_basics}

\section{什么是 Linux}
\label{sec:what_is_linux}

Linux 是一个开源的类 Unix 操作系统内核，由 Linus Torvalds 于1991年创建。

Linux 的特点：
\begin{itemize}
    \item \textbf{开源}：源代码开放，自由使用和修改
    \item \textbf{稳定}：稳定性高，适合服务器使用
    \item \textbf{安全}：内置安全机制
    \item \textbf{多用户}：支持多用户同时使用
    \item \textbf{多任务}：支持多任务并行执行
\end{itemize}

\section{Linux 发行版}
\label{sec:linux_distributions}

常见的 Linux 发行版：
\begin{table}[htbp]
    \centering
    \caption{Linux 发行版}
    \begin{tabular}{|c|p{8cm}|}
        \hline
        \textbf{发行版} & \textbf{说明} \\
        \hline
        CentOS/Rocky Linux & 企业级服务器系统 \\
        \hline
        Ubuntu & 桌面和服务器都适用 \\
        \hline
        Debian & 稳定可靠的发行版 \\
        \hline
        Fedora & 创新型发行版 \\
        \hline
        openSUSE & 企业级发行版 \\
        \hline
        Arch Linux & 滚动更新，面向高级用户 \\
        \hline
    \end{tabular}
\end{table}

\section{Linux 内核版本}
\label{sec:kernel_version}

内核版本格式：主版本号.次版本号.修订号

查看内核版本：
\begin{verbatim}
uname -r
cat /proc/version
\end{verbatim}

\chapter{命令行基础}
\label{chap:command_line}

\section{Shell 简介}
\label{sec:shell_intro}

Shell 是用户与内核之间的接口：
\begin{table}[htbp]
    \centering
    \caption{Shell 类型}
    \begin{tabular}{|c|p{8cm}|}
        \hline
        \textbf{Shell} & \textbf{说明} \\
        \hline
        Bash & 最常用的 Shell \\
        \hline
        Zsh & 功能丰富的 Shell \\
        \hline
        Fish & 用户友好的 Shell \\
        \hline
    \end{tabular}
\end{table}

查看当前 Shell：
\begin{verbatim}
echo $SHELL
\end{verbatim}

切换 Shell：
\begin{verbatim}
chsh -s /bin/zsh
\end{verbatim}

\section{文件系统导航}
\label{sec:file_navigation}

常见目录：
\begin{table}[htbp]
    \centering
    \caption{Linux 目录结构}
    \begin{tabular}{|c|p{8cm}|}
        \hline
        \textbf{目录} & \textbf{用途} \\
        \hline
        / & 根目录 \\
        \hline
        /home & 用户主目录 \\
        \hline
        /etc & 系统配置文件 \\
        \hline
        /var & 可变数据（日志等） \\
        \hline
        /usr & 用户程序和数据 \\
        \hline
        /bin & 二进制可执行文件 \\
        \hline
        /sbin & 系统二进制文件 \\
        \hline
        /tmp & 临时文件 \\
        \hline
    \end{tabular}
\end{table}

常用命令：
\begin{verbatim}
pwd                 # 显示当前目录
cd /path/to/dir     # 切换目录
ls                  # 列出文件
ls -la              # 列出详细信息
mkdir dirname       # 创建目录
rmdir dirname       # 删除目录
\end{verbatim}

\section{文件操作}
\label{sec:file_operations}

常用文件操作命令：
\begin{verbatim}
touch filename      # 创建文件
cat filename        # 查看文件内容
more filename       # 分页查看
less filename       # 交互式查看
cp src dest         # 复制文件
mv src dest         # 移动文件
rm filename         # 删除文件
rm -rf dirname      # 递归删除目录
\end{verbatim}

查看文件头部和尾部：
\begin{verbatim}
head -n 10 filename
tail -n 10 filename
tail -f filename    # 实时查看
\end{verbatim}

\chapter{用户与权限}
\label{chap:users_and_permissions}

\section{用户管理}
\label{sec:user_management}

用户类型：
\begin{itemize}
    \item \textbf{根用户}：UID 0，拥有全部权限
    \item \textbf{普通用户}：UID 1000+
    \item \textbf{系统用户}：UID 1-999
\end{itemize}

用户管理命令：
\begin{verbatim}
useradd username    # 创建用户
userdel username    # 删除用户
usermod -aG group username  # 添加到组
passwd username     # 设置密码
id username         # 查看用户信息
whoami              # 当前用户
\end{verbatim}

切换用户：
\begin{verbatim}
su - username       # 切换用户
sudo command        # 以 root 身份执行
\end{verbatim}

\section{权限管理}
\label{sec:permission_management}

文件权限：
\begin{itemize}
    \item \textbf{r}：读取权限 (4)
    \item \textbf{w}：写入权限 (2)
    \item \textbf{x}：执行权限 (1)
\end{itemize}

权限设置：
\begin{verbatim}
chmod 755 filename  # 设置权限
chmod +x filename   # 添加执行权限
chown user:group filename  # 更改所有者
chgrp group filename       # 更改组
\end{verbatim}

\chapter{进程管理}
\label{chap:process_management}

\section{进程基础}
\label{sec:process_basics}

进程是正在运行的程序实例。

查看进程：
\begin{verbatim}
ps aux              # 列出所有进程
ps aux | grep name  # 查找进程
top                 # 实时进程监控
htop                # 交互式进程监控
\end{verbatim}

进程状态：
\begin{table}[htbp]
    \centering
    \caption{进程状态}
    \begin{tabular}{|c|p{8cm}|}
        \hline
        \textbf{状态} & \textbf{说明} \\
        \hline
        R & 运行中 \\
        \hline
        S & 睡眠中 \\
        \hline
        D & 不可中断睡眠 \\
        \hline
        Z & 僵尸进程 \\
        \hline
        T & 停止 \\
        \hline
    \end{tabular}
\end{table}

\section{进程控制}
\label{sec:process_control}

进程控制命令：
\begin{verbatim}
kill PID            # 终止进程
kill -9 PID         # 强制终止
killall name        # 按名称终止
pkill name          # 按名称终止
\end{verbatim}

前后台进程：
\begin{verbatim}
command &           # 后台运行
jobs                # 查看后台任务
fg %1               # 切换到前台
bg %1               # 切换到后台
\end{verbatim}

\chapter{软件包管理}
\label{chap:package_management}

\section{RPM 包管理}
\label{sec:rpm_management}

RPM（Red Hat Package Manager）：

常用命令（CentOS/Rocky Linux）：
\begin{verbatim}
dnf install package     # 安装包
dnf remove package      # 删除包
dnf update package      # 更新包
dnf list installed      # 已安装包
dnf search package      # 搜索包
dnf info package        # 包信息
dnf clean all           # 清理缓存
\end{verbatim}

安装本地包：
\begin{verbatim}
dnf install ./package.rpm
rpm -ivh package.rpm
\end{verbatim}

\section{DEB 包管理}
\label{sec:deb_management}

DEB（Debian Package）：

常用命令（Ubuntu/Debian）：
\begin{verbatim}
apt update             # 更新源
apt upgrade            # 升级包
apt install package    # 安装包
apt remove package     # 删除包
apt autoremove         # 删除无用包
apt search package     # 搜索包
apt show package       # 包信息
\end{verbatim}

\chapter{系统配置}
\label{chap:system_config}

\section{网络配置}
\label{sec:network_config}

配置网络接口：
\begin{verbatim}
nmcli connection show              # 查看连接
nmcli connection add type ethernet con-name eth0 ifname eth0
nmcli connection modify eth0 ipv4.addresses 192.168.1.100/24
nmcli connection modify eth0 ipv4.gateway 192.168.1.1
nmcli connection modify eth0 ipv4.dns 8.8.8.8
nmcli connection modify eth0 ipv4.method manual
nmcli connection up eth0
\end{verbatim}

查看网络信息：
\begin{verbatim}
ip addr show
ip route show
cat /etc/resolv.conf
\end{verbatim}

\section{时间配置}
\label{sec:time_config}

查看时间：
\begin{verbatim}
date
timedatectl
\end{verbatim}

设置时区：
\begin{verbatim}
timedatectl set-timezone Asia/Shanghai
\end{verbatim}

同步时间：
\begin{verbatim}
dnf install -y chrony
systemctl enable --now chronyd
chronyc sources
\end{verbatim}

\section{防火墙配置}
\label{sec:firewall_config}

firewalld 命令：
\begin{verbatim}
firewall-cmd --list-all          # 查看规则
firewall-cmd --add-service=http --permanent
firewall-cmd --add-port=8080/tcp --permanent
firewall-cmd --remove-service=http --permanent
firewall-cmd --reload            # 重载规则
\end{verbatim}

\chapter{服务管理}
\label{chap:service_management}

\section{systemd}
\label{sec:systemd}

systemd 是 Linux 系统管理器：

常用命令：
\begin{verbatim}
systemctl status service    # 查看服务状态
systemctl start service     # 启动服务
systemctl stop service      # 停止服务
systemctl restart service   # 重启服务
systemctl enable service    # 开机自启
systemctl disable service   # 禁用开机自启
systemctl list-units --type=service  # 服务列表
\end{verbatim}

查看系统日志：
\begin{verbatim}
journalctl -u service
journalctl -f              # 实时日志
journalctl --since today   # 今天的日志
\end{verbatim}

\section{自定义服务}
\label{sec:custom_service}

创建自定义服务：
\begin{verbatim}
cat > /etc/systemd/system/myapp.service <<EOF
[Unit]
Description=My Application
After=network.target

[Service]
Type=simple
ExecStart=/usr/bin/myapp
Restart=always
User=myuser

[Install]
WantedBy=multi-user.target
EOF

systemctl daemon-reload
systemctl enable --now myapp
\end{verbatim}

\chapter{存储管理}
\label{chap:storage}

\section{磁盘管理}
\label{sec:disk_management}

查看磁盘信息：
\begin{verbatim}
lsblk
fdisk -l
df -h
du -sh /path
\end{verbatim}

分区操作：
\begin{verbatim}
fdisk /dev/sdb
n    # 创建新分区
p    # 主分区
w    # 保存
\end{verbatim}

格式化分区：
\begin{verbatim}
mkfs.xfs /dev/sdb1
mkfs.ext4 /dev/sdb1
\end{verbatim}

挂载分区：
\begin{verbatim}
mount /dev/sdb1 /data
echo '/dev/sdb1 /data xfs defaults 0 0' >> /etc/fstab
\end{verbatim}

\section{LVM 管理}
\label{sec:lvm_management}

LVM 命令：
\begin{verbatim}
pvcreate /dev/sdb1
vgcreate vgdata /dev/sdb1
lvcreate -L 50G -n lvdata vgdata
mkfs.xfs /dev/vgdata/lvdata
mount /dev/vgdata/lvdata /data
\end{verbatim}

扩展卷：
\begin{verbatim}
lvextend -L +20G /dev/vgdata/lvdata
xfs_growfs /data
\end{verbatim}

\chapter{性能监控与优化}
\label{chap:performance}

\section{性能监控工具}
\label{sec:monitoring_tools}

常用监控工具：
\begin{table}[htbp]
    \centering
    \caption{性能监控工具}
    \begin{tabular}{|c|p{8cm}|}
        \hline
        \textbf{工具} & \textbf{说明} \\
        \hline
        top & 进程监控 \\
        \hline
        htop & 交互式进程监控 \\
        \hline
        iostat & 磁盘 I/O \\
        \hline
        vmstat & 内存和 CPU \\
        \hline
        free & 内存使用 \\
        \hline
        mpstat & CPU 统计 \\
        \hline
    \end{tabular}
\end{table}

查看内存：
\begin{verbatim}
free -h
cat /proc/meminfo
\end{verbatim}

查看 CPU：
\begin{verbatim}
cat /proc/cpuinfo
nproc
\end{verbatim}

\section{性能优化}
\label{sec:optimization}

性能优化策略：
\begin{itemize}
    \item \textbf{内存优化}：调整 swappiness
    \item \textbf{磁盘优化}：使用 SSD、RAID
    \item \textbf{网络优化}：调整 TCP 参数
    \item \textbf{进程优化}：调整优先级
\end{itemize}

调整进程优先级：
\begin{verbatim}
nice -n 10 command    # 低优先级
renice -n 10 PID      # 调整优先级
\end{verbatim}

\backmatter

\end{document}